\title{LongRoPE: Extending LLM Context Window Beyond 2 Million Tokens}

\begin{abstract}
Expanding the context window is an advantageous characteristic of large language models (LLMs). Nevertheless, existing long context windows generally reach up to 128k tokens, primarily constrained by expensive fine-tuning processes, limited availability of lengthy training samples, and the risk of problematic values arising from newly introduced token positions.

This study presents LongRoPE, a novel approach that, for the first time, expands the context window of pre-trained LLMs to an outstanding \textbf{2048k} tokens. Remarkably, this is accomplished with as few as 1k fine-tuning steps on training sequences of up to 256k tokens, all while preserving the model’s effectiveness within its original, shorter context window. The method is underpinned by three central innovations: (i) through an efficient search, we uncover and utilize two distinct non-uniformities in positional interpolation, thereby enabling superior fine-tuning initialization and facilitating an 8$\times$ context window increase even without fine-tuning; (ii) we propose a stepwise expansion technique, in which a model is first fine-tuned for 256k token sequences, followed by an additional round of positional interpolation on the newly fine-tuned model, culminating in support for a 2048k token window; (iii) we further adapt LongRoPE on sequences of 8k length to restore optimal performance for short contexts. Comprehensive empirical evaluations on LLaMA2 and Mistral, spanning multiple tasks, validate the proposed approach. Models enhanced with LongRoPE require only minimal changes to the positional embedding layer and maintain the original network architecture, thus retaining compatibility with most established optimization strategies. The implementation will be released at \texttt{https://github.com/microsoft/LongRoPE}

\section{Introduction}

Although Large Language Models (LLMs) have demonstrated impressive capabilities across a range of tasks~\cite{gpt4,llama2}, they are constrained by a relatively small context window, such as the 4096-token maximum in LLaMA2~\cite{llama2}. When processing text that extends beyond this limit, the models exhibit reduced performance, largely because they encounter positions during inference that were not encountered during training. This limitation becomes particularly problematic in applications where extensive in-context examples are required~\cite{Huang2023FewerIM}, or in the operation of LLM-powered agents~\cite{park2023generative,madaan2023selfrefine}.

Recent studies indicate that the context window of a pre-trained LLM can be stretched to approximately 128k tokens through fine-tuning with extended-length documents~\cite{longlora,pi,yarn,zhang2024soaring,liu2023scaling}. However, three primary challenges hinder further expansion. First, incorporating uninitialized position indices results in numerous anomalous values, which triggers out-of-distribution complications and impedes the convergence of fine-tuning~\cite{pi}. This complication becomes especially severe when upgrading from 4k to $>$1000k positions, as the process introduces over 90\% new, untrained indices. Second, fine-tuning generally depends on the availability of texts with matching lengths, but current datasets rarely feature examples beyond 1000k tokens. Furthermore, the process of training on these ultra-long texts is computationally intensive, necessitating excessive durations and substantial GPU resources. Third, as the context window is pushed to extreme lengths, attention mechanisms distribute across an ever-growing array of token positions, which diminishes performance for shorter, original contexts~\cite{pi}.

To address the initial challenge, a common method involves interpolating RoPE positional embeddings~\cite{rope,pi}. This technique scales down newly introduced position indices so that they fit within the pretrained interval, as illustrated in Fig.\ref{fig:overview}. Position Interpolation (PI)~\cite{pi} implements a linear adjustment of RoPE's rotary angle values based on the ratio of extension. NTK~\cite{ntk,dynamicntk} proposes that interpolation and extrapolation should vary across RoPE dimensions. YaRN~\cite{yarn} segments the RoPE dimensions according to three frequency tiers, employing different schemes—extrapolation, NTK, or linear interpolation—on each group. Despite these innovations, positional embeddings within the Transformer framework contain \emph{complex, non-uniform information entropy} that is not fully exploited by current methodologies. As a result, this nuanced entropy is lost, which constrains the expansion of the context window.

Section~\ref{sec:analysis} empirically demonstrates two principal results: \textbf{(1)} For positional interpolation to be effective, it is necessary to account for two types of non-uniformity: discrepancies among RoPE dimensions and among token positions. Minimal interpolation is required for lower RoPE dimensions as well as for early token positions, although the ideal configuration is contingent on the intended extended sequence length. \textbf{(2)} By integrating these forms of non-uniformity into the positional interpolation strategy, it becomes possible to better preserve crucial information present in the initial RoPE—especially in key dimensions and at salient token positions. This approach mitigates information loss stemming from positional interpolation, thereby leading to improved initial weights for subsequent fine-tuning. Additionally, it supports up to 8$\times$ context length extension in cases where fine-tuning is not employed.

Building on these results, we propose \textbf{LongRoPE}, a robust approach capable of increasing the LLM context window to over 2 \emph{million} tokens. The design of LongRoPE incorporates three primary advancements. Firstly, it leverages multidimensional disparities in positional interpolation to their fullest extent. Specifically, LongRoPE determines optimal rescale coefficients for the rotational parameters within RoPE for each respective dimension, taking token location into account. Given that the complexity of rescale factor identification grows exponentially as the desired extension ratio rises, LongRoPE implements an evolutionary search algorithm supplemented by two optimization strategies, enhancing the efficiency of the search process. An illustration of the resulting rescaled RoPE produced by this method is presented in Fig.~\ref{fig:overview}.

Subsequently, {LongRoPE} utilizes a resource-efficient, incremental extension methodology to realize a 2048k context window. This is accomplished without performing direct fine-tuning on extremely lengthy texts, which are rare and not easily obtainable. The approach initiates by identifying a suitable 256k sequence length using the pre-trained LLM, followed by fine-tuning at this length. Due to the non-uniform positional interpolation, which enables an 8$\times$ extension even in the absence of fine-tuning, a second phase of searching for optimal RoPE rescale factors is conducted on the extended, fine-tuned LLM. Through this process, the model attains the 2048k context window capability for LLaMA2 and Mistral~\cite{mistral}.

In conclusion, {LongRoPE} addresses potential declines in performance on standard (shorter) context lengths by further optimizing the RoPE rescale factors in the extended LLM. Analogous to the upscaling process from 256k to 2048k, we apply our search method to scale down to 4k and 8k context windows on the LLM fine-tuned at 256k, aiming to minimize positional interpolation. When inference sequences are shorter than 8k, we apply the rescale factors identified by our search procedure to update RoPE accordingly.

A comprehensive set of experiments conducted on multiple LLMs and a range of long-context tasks highlights the strong performance of our approach. Our results indicate that LongRoPE preserves low perplexity even as the evaluation length scales from 4k up to 2048k tokens, attains passkey retrieval accuracy exceeding 90\%, and maintains competitive accuracy on established benchmarks with a 4096-token context window. Importantly, LongRoPE is compatible with all LLMs employing RoPE embedding. We plan to make our code and the LongRoPE-2048k models publicly available.

\section{Non-uniformity in Positional Interpolation}
\label{sec:analysis}

\subsection{Preliminary}

Transformer models require explicit positional information, often in the form of position embedding, to represent the order of input tokens. Our work focuses on the RoPE~\cite{rope} position embedding, which is widely used in recent  LLMs. For a token at position index $n$, its corresponding RoPE encoding can be simplified as follows:

\begin{equation}
		\fontsize{7.8}{7.8} \selectfont
	\label{eq:rope}
	\left[ cos(n\theta_0), sin(n\theta_0), cos(n\theta_1), \cdot\cdot\cdot, cos(n\theta_{d/2-1}), sin(n\theta_{d/2-1}) \right]   
\end{equation}

Here, $d$ denotes the dimension of the embedding, $n\theta_i$ specifies the rotary angle associated with the token positioned at $n$, and $\theta_i=\theta^{-2i/d}$ characterizes the rotation frequencies. RoPE typically utilizes a default base $\theta$ of 10000.

\textbf{\emph{Context window extension ratio $s$} and positional interpolation}. We introduce $s$ as the proportion of the extended context length $L'$ compared to the original context length $L$, defined as $s = \frac{L'}{L}$.

To extend the context window from $L$ to $L'$, current positional interpolation methods suggest downscaling rotation frequencies $\theta_i$ based on extension ratio $s$. Let $\beta=\theta^{2/d}$, and $\lambda$ denote the actual rescale factor related to $s$, we   unify these positional interpolation methods as follows:

\begin{equation}	
	\fontsize{7.5}{7.5} \selectfont
	\label{eq:unified_rope}
	\begin{aligned}
		\left[cos\left(\frac{n}{\lambda(\beta)^0}\right), sin \left(\frac{n}{\lambda(\beta)^0}\right), cos\left(\frac{n}{\lambda(\beta)^1}\right), \cdot\cdot\cdot,  sin \left(\frac{n}{\lambda(\beta)^{d/2-1}}\right) \right]   
	\end{aligned}
\end{equation}

\textbf{Linear positional interpolation (PI)}. The PI method~\cite{pi} involves applying a linear scaling to position indices, restricting them within the original training sequence length. When extending to a longer context with a ratio $s$, each rotation angle corresponding to token positions is multiplied by $\lambda=s$ across all RoPE dimensions. This approach compresses the positional encoding, resulting in densely packed position signals that impede the model's capacity to differentiate tokens that are near each other. As a consequence, PI generally yields suboptimal results when the extension ratio is high.

\textbf{NTK-based interpolation and extrapolation}. Works such as ~\cite{ntk,dynamicntk} approach RoPE through the lens of information encoding, leveraging Neural Tangent Kernel (NTK) theory~\cite{jacot2018neural,tancik2020fourier}. To address the problem of position crowding in PI, these studies propose spreading the interpolation influence among the various RoPE dimensions. Specifically, lower (high-frequency) dimensions are scaled less, while higher (low-frequency) dimensions are scaled more, facilitating both positional interpolation and extrapolation through $\lambda=s^{i}$. The enhanced dynamic NTK~\cite{dynamicntk} adapts the extension ratio for each position according to the given sequence length. In contrast to PI, which relies on fine-tuning, NTK-aware strategies permit context window extension even in settings where fine-tuning is not performed, albeit typically allowing for a maximum expansion factor of 4$\times$.

\textbf{YaRN}~\cite{yarn} presents an advance in enhancing the efficacy of positional interpolation. The method segments RoPE dimensions according to frequency, assigning each segment a tailored interpolation approach. Specifically, dimensions with high frequencies are processed using extrapolation ($\lambda$=1), whereas those with low frequencies utilize linear interpolation (PI). For intermediate frequencies, the NTK interpolation method is applied. Central to YaRN’s methodology is the classification of RoPE dimensions into these groups, a procedure currently reliant on empirical human judgment, which could limit optimal performance when adapting to novel LLMs.

\subsection{Study on Non-uniform Positional Interpolation}

Drawing inspiration from NTK and YaRN, we observe that their performance improvements stem from incorporating non-linearity, particularly through accounting for varying frequencies in RoPE dimensions to tailor interpolation and extrapolation. Yet, prevailing approaches to non-linearity depend significantly on manually crafted heuristics. This observation leads to two pertinent inquiries: (1) Does existing positional interpolation achieve optimality? (2) Might there exist as-yet-uninvestigated forms of non-linearity?

In order to address these questions, we employ evolution search (refer to Sec.~\ref{sec:method}) to identify improved non-uniform positional interpolation strategies for LLaMA2-7B. This search process is directed by evaluating perplexity on 5 randomly selected passages from the PG19~\cite{pg19} validation dataset. Our experiments yield several significant observations, detailed below.

\textbf{Finding 1}: \textit{RoPE dimensions exhibit substantial non-uniformities, which are not effectively handled by current positional interpolation methods.}

We optimize $\lambda$ for each RoPE dimension as specified in Eq.~\ref{eq:unified_rope}. In Table~\ref{tbl:exp1}, we report the perplexity scores of LLaMA2-7B on PG19 and Proof-pile~\cite{proof-pile} test data, across different methods, all evaluated without any additional fine-tuning. Our search strategy yields marked perplexity reductions, indicating that commonly used linear (PI) and irregular (Dynamic-NTK and YaRN) interpolation strategies are not ideal. Interestingly, on the PG19 dataset, YaRN yields higher perplexity compared to PI and NTK, due to its inability to extend to the desired context length for language models that have not undergone further fine-tuning. Specifically, for an 8k context length, we observe a pronounced increase in perplexity for YaRN beyond 7k tokens.

In our exploration, the rescaling coefficients $\lambda$ in Eq.~\ref{eq:unified_rope} are no longer uniform; they differ from the constant scale $s$ applied in PI, as well as the formulas used in NTK and the group-specific adjustments in YaRN. Adopting these varied rescale factors leads to notable improvements in LLaMA2's language modeling, as demonstrated by reductions in perplexity at 8k and 16k context lengths, all without additional fine-tuning. This advantage stems from the revised positional embedding's capacity to better maintain the characteristics of the original RoPE—particularly in crucial dimensions—thereby alleviating the model's challenge in accurately differentiating between nearby token positions.

\textbf{Finding 2}: \textit{RoPE for the initial tokens in the input sequence should be extrapolated with less interpolation.} 

We posit that for the first $\hat{n}$ tokens in the input, RoPE should employ minimal interpolation. This is due to the high attention weights assigned to these early tokens, underscoring their importance within the attention mechanism, as previously reported in Streaming LLM~\cite{streamingllm} and LM-Infinite~\cite{han2023lminfinite}. To empirically test this, we increase the context window to 8k and 16k via PI and NTK methods, while deliberately withholding interpolation for the initial $\hat n$ tokens ($0, 2, \ldots, 256$). Setting $\hat n=0$ represents the baseline configurations of PI and NTK. Our findings are presented in Table~\ref{tbl:tokenposition}, which reveal: \textbf{(1)} excluding the earliest tokens from position interpolation leads to superior performance; and \textbf{(2)} the best value for $\hat n$ varies based on how far the context window is extended.

\textbf{Finding 3}: \textit{Non-uniform positional interpolation effectively extends LLM context window in both fine-tuning and non-fine-tuning settings.} 

Our experiments demonstrate that the use of our optimized non-uniform position interpolation leads to notable gains in context extension at 8k and 16k, even without additional fine-tuning. However, extending to larger context windows necessitates fine-tuning. To address this, we conduct fine-tuning of LLaMA2-7B employing our searched RoPE configuration for a context window of 64k tokens (refer to Appendix for further details). Results presented in Table~\ref{tbl:finetune} indicate that our approach achieves markedly better performance than PI and YaRN, both in the pre- and post-fine-tuning phases for LLaMA2-7B. These results stem from our method’s more precise application of non-uniform positional interpolation, which reduces information loss and supplies an advantageous initialization prior to fine-tuning.

\textbf{Summary}. This work identifies two core non-uniformities: differences in RoPE dimensions and in token positions. Effectively leveraging these non-uniform patterns in positional interpolation produces substantial improvements in extending LLM context length.

\section{LongRoPE}

\label{sec:method}
Building on these observations, we propose LongRoPE. This approach incorporates an optimized search algorithm designed to leverage both sources of non-uniformity, which enables scaling the context window of LLMs to exceed 2 million tokens.

\subsection{Problem Formulation}

These two sources of non-uniformity result in an expansive array of possible solutions and add challenges to the optimization process. To mitigate these issues, we reinterpret the multidimensional non-uniform position interpolation optimization task as a search problem.

For a LLM targeting a  context window size of $L'$ and lengthy input documents $\mathbf{X}$, where each $\mathbf{x}\in\mathbf{X}$ surpasses $L'$ in token length, we denote the original rotary angle of the $i^{th}$ dimension in RoPE embedding at token position $n$ as  $\frac{n}{\beta^i}$. The optimization problem is then formulated as follows:

\begin{equation}
\begin{gathered}
		\label{eq:problem}

\underset{\mathbf{x}\in\mathbf{X}; \, |\mathbf{x}|\ge L'}{\operatorname{arg\,min}}\, \mathcal{L}\,\left( \text{LLM} (\text{RoPE}, \mathbf{X}) \right), \text{where}
\\
\scalebox{0.93}{
$\underset{\begin{subarray}{l} i=0, \cdot\cdot, \frac{d}{2}-1;\\ n\in [0, |\mathbf{x}|); \end{subarray}}{\text{RoPE}_{(n)}} = {\left[\cdot\cdot, \cos\left( \mathbb{I}(\hat{\lambda}_{i}, \hat{n}) \cdot \frac{n}{\beta^i} \right), \sin\left( \mathbb{I}(\hat{\lambda}_{i}, \hat{n}) \cdot \frac{n}{\beta^i} \right), \cdot\cdot\right]}$} \\
\text{where}\; \mathbb{I}(\hat{\lambda}_{i},\hat{n}) = \begin{cases}
1 &\mbox{ if $n < \hat{n}$} \\
{\frac{1}{\lambda}_{i}} &\mbox{ if $n \geq \hat{n}$}
\end{cases}
\end{gathered}
\end{equation}

In this approach, we define a group of scaling coefficients, $\mathbb{I}(\hat{\lambda}_{i},\hat{n})$, to address both types of non-uniformities. Here, $\hat{\lambda}_{i}$ reflects the variability among RoPE dimensions, while $\hat{n}$ captures the irregularity across token positions. The scaling function $\mathbb{I}(\hat{\lambda}_{i},\hat{n})$ adjusts the rotation angle corresponding to the $i^{th}$ dimension of RoPE; specifically, $\hat{\lambda}_{i}$ is the factor used for scaling, and $\hat{n}$ sets the threshold for token positions. For the first $\hat{n}-1$ tokens, no scaling is applied, so the standard RoPE rotation angle $\frac{n}{\beta^i}$ remains unchanged. When token positions satisfy $n \geq \hat{n}$, the rescale factor is introduced.

Assuming a desired context window length of $L'$, the task involves determining the best set of rescale factors ($\mathbb{I}(\hat{\lambda}_{0},\hat{n})$, $\mathbb{I}(\hat{\lambda}_{1},\hat{n})$, ..., $\mathbb{I}(\hat{\lambda}_{i},\hat{n})$, ...) across the $1^{st}$ to $d^{th}$ dimensions of RoPE. Through this process, the target $\text{LLM}$ utilizing the adjusted $\text{RoPE}$ can minimize the next token prediction loss, $\mathcal{L}$ (quantified as perplexity), over the given input sequences $\mathbf{X}$ of length $L'$.

\subsection{Searching the Non-uniform Position Interpolation}

In addressing the challenge presented in Eq.~\ref{eq:problem}, we present a straightforward and potent approach that systematically identifies the best RoPE rescale factors. This strategy harnesses the full potential of multidimensional disparities within position embeddings.

\textbf{Search space}. We construct an extensive search space that accommodates the two types of non-uniformity. Table~\ref{tbl:searchspace} provides a detailed overview of this search space. Our approach enables independent tuning of the rescale factor for each dimension within the RoPE embedding. To simplify the configuration of the search space, we focus on optimizing $\lambda_i$ and $\hat{n}$, rather than directly adjusting $\mathbb{I}(\hat{\lambda}_{i},\hat{n})$, with the relationship $\hat{\lambda}_{i}=1/\lambda_i$. As indicated in Table~\ref{tbl:searchspace}, the value for $\lambda_i$ is selectable from 1.0 (representing simple extrapolation) up to $s\times1.25$ (which represents greater interpolation than PI), incrementing by 0.01. Here, $s$ denotes the intended context window extension ratio.

The parameter $\hat{n}$ determines how many of the earliest token positions are kept with their original position encoding (that is, employing the unaltered RoPE embedding instead of interpolation). In our experiments, $\hat{n}$ is selected from the set \{0, 1, 2, 4, 8, 12, 16, 20, 24, 28, 32, 64, 128, 256\}. Setting $\hat{n}=0$ implies that every token position applies the learned rescaling factors.

\textbf{Evolution-based search}. The search space detailed in Table~\ref{tbl:searchspace} encompasses a vast array of positional interpolation options, making the process of effective search particularly demanding. As an illustration, an extension factor of $s=4\times$ produces $400^{128/2}\times14$ = $4\times10^{167}$ possible combinations. The size of this space grows exponentially as the extension ratio increases. To overcome these complexities, we employ an evolution search method~\cite{guo2020single} and implement two optimization strategies that significantly enhance search performance. Algorithm~\ref{alg:search} outlines the complete search methodology.

\textit{Optimized initial population generation}. Rather than relying solely on random assignment for the initial population of $P$ rescale factors, we explicitly include the three RoPE rescale factors associated with PI, NTK, and YaRN as predetermined individuals within the initial population. The other $P$-3 members are created by subjecting these three rescale factors to random mutations with probability $p$.

\textit{Monotonically non-decreasing constraint}. Once the initial population is formed, LLM perplexity values are assessed for every candidate. This involves using the appropriate RoPE rescale factors on the target LLM and evaluating perplexity for the input $\mathbf{X}$. The best $k$ individuals are chosen as parent candidates for evolutionary steps. Nevertheless, due to the expansive nature of the search space, uninformed mutation and crossover operations may frequently sample suboptimal regions, resulting in superfluous perplexity evaluations. Such inefficiency is further amplified as $L'$ increases, given the heavy computational load associated with each perplexity inference.

To mitigate this issue, we enforce a monotonic non-decreasing constraint on the sequence of sampled RoPE rescaled factors such that $\lambda_i \leq \lambda_{i+1}$. Only the RoPE configurations adhering to this condition are utilized during perplexity assessment within the LLM, thereby considerably diminishing computational search overhead. In detail, the constraint mandates that $\lambda_i$ should monotonically rise with increasing RoPE dimension indices ($i=0,\ldots,63$). This requirement stems from NTK theory~\cite{jacot2018neural,tancik2020fourier,ntk}, which posits that lower-dimensional components—corresponding to higher frequencies—necessitate minimal interpolation (yielding smaller $\lambda_i$), whereas higher-dimensional components—with lower frequencies—are able to accommodate greater interpolation (resulting in larger $\lambda_i$).

\begin{algorithm}[t]
	\small
	\caption{The search algorithm}
	\textbf{Input:} target LLM, input samples $\mathbf{X}$, population size $P$, mutation size $N_1$, crossover size $N_2$, maximum iterations $\mathcal{T}$, mutation probability $p$\\

	\begin{algorithmic}[1]
		\label{alg:search}
		\STATE Initialize Top-k as $\phi$;
		\STATE Generate initial population $\text{P}_0$ using \textit{Initialize\_population\_with\_optimization} ($P$, $\mathbf{X}$, $p$);
		\FOR{$i$ from $1$ to $\mathcal{T}$}
		\STATE Evaluate perplexity via \textit{Compute\_perplexity} (LLM, $\text{P}_{i-1}$, $\mathbf{X}$);
		\STATE Refresh Top-k with \textit{Update\_Topk} (Top-k, $\text{P}_{i-1}$);
		\STATE Create mutated population $\text{P}_{mutation}$ by applying \textit{Mutation\_with\_mono\_constraint} (Top-k, $N_1$, $p$);
		\STATE Produce crossover population $\text{P}_{crossover}$ using \textit{Crossover\_with\_mono\_constraint} (Top-k, $N_2$);
		\STATE Set $\text{P}_i$ to the union of $\text{P}_{mutation}$, $\text{P}_{crossover}$, and Top-k;
		\ENDFOR
		\STATE Output the candidate from Top-k with the minimum perplexity;
	\end{algorithmic}
\end{algorithm}

\textbf{8$\times$ extension achieved without fine-tuning}. The evolutionary search approach employed here selects optimal non-uniform RoPE rescale factors, ensuring essential dimensions and positions are maintained to reduce information loss from interpolation. Figure~\ref{fig:non-ft} illustrates that this method successfully expands LLaMA2's context window from 4k to 32k with no fine-tuning required. By comparison, methods like PI, non-uniform NTK, and YaRN exhibit a sharp increase in perplexity when the context window is extended beyond 2$\times$.

\subsection{Extending LLM Context Window to 2048K}
\label{sec:extendto2048k}

\textbf{Progressive extension to 2048k}. In this section, we propose an approach for increasing the context window of pre-trained LLMs from the standard 4k tokens up to 2048k tokens. Our technique, which employs non-uniform positional interpolation, allows for an 8$\times$ increase in context length without any requirement for fine-tuning. However, when scaling further to, for example, 512$\times$ extensions, fine-tuning becomes indispensable. One approach involves identifying suitable RoPE rescale parameters for the 2048k context size and then performing fine-tuning. Despite this, such an approach is hindered by the substantial computational and resource demands of training at this scale. Furthermore, our empirical observations indicate that effectively fine-tuning LLMs for such large context window expansions is particularly difficult (refer to Appendix for details).

Fortunately, LongRoPE demonstrates efficacy with both the baseline and fine-tuned versions of extended LLM. Consequently, we propose a streamlined, incremental approach that reaches the desired 2048k context window utilizing only 1k fine-tuning steps conducted within a training sequence of 256k length.

$\Diamond$ \textit{LongRoPE search for expanding pre-trained LLM to 256k}. Using LLaMA2 for illustration, we perform a search targeting context windows of 128k and 256k. The resulting extension factors are 32$\times$ and 64$\times$ at this phase, respectively.

$\Diamond$ \textit{Fine-tuning to 256k}. Next, the pre-trained LLM undergoes fine-tuning to expand the context window to 256k. To accomplish this, we initially fine-tune LLaMA2 with RoPE rescaled factors adjusted for 128k over 400 steps. After completing this phase, we update the RoPE rescaled factors to support 256k at the checkpoint and proceed with a further 600 fine-tuning steps. This incremental approach demonstrates greater efficiency compared to a direct fine-tuning for 256k context length.

$\Diamond$ \textit{Scaling up fine-tuned extended LLM to 2048k via LongRoPE search}. In the final step, we apply an additional search phase to the previously fine-tuned LLM with a 256k context length. As a result, the context window can be expanded to a remarkable 2048k length, achieved without the need for extra fine-tuning. This process yields a total extension factor of $512\times$.

\textbf{Restoring performance in shorter context windows}.  Expanding the context window to an extensive 2048k length results in a noticeable decline in model performance inside the initial window size. This drawback of positional interpolation~\cite{pi} arises because position embeddings in the early segment are compressed into a smaller region in the higher-dimensional space, which can hamper the quality of language modeling. When the context window is scaled by a factor of 512, the position embeddings associated with the original 4k window become densely packed.

To address this issue, an additional evolutionary search is conducted on the expanded LLM in order to tune RoPE rescale factors specifically for shorter context lengths such as 4k and 8k. The upper bound for the searched $\lambda$ is lowered, reflecting the reduced need for positional interpolation with these smaller lengths. When inferring, the LLM adaptively modifies the relevant RoPE rescale factors in response to context.

\section{LongRoPE}

\label{sec:method}
Building upon these observations, we propose LongRoPE, which initially employs an effective search algorithm designed to leverage both forms of non-uniformity. This approach enables the extension of the LLM context window to exceed 2 million tokens.

\subsection{Problem Formulation}

These two types of non-uniformity expand the solution space significantly and complicate the optimization process. To manage this challenge, we reformulate the multidimensional non-uniform position interpolation optimization as a search problem.

For a LLM targeting a  context window size of $L'$ and lengthy input documents $\mathbf{X}$, where each $\mathbf{x}\in\mathbf{X}$ surpasses $L'$ in token length, we denote the original rotary angle of the $i^{th}$ dimension in RoPE embedding at token position $n$ as  $\frac{n}{\beta^i}$. The optimization problem is then formulated as follows:

\begin{equation}
\begin{gathered}
		\label{eq:problem}

\underset {\mathbf{x}\in\mathbf{X}; \, |\mathbf{x}|\ge L'}{ \operatorname {arg\,min} } \,  \mathcal{L}\left( \text{LLM} (\text{RoPE}, \mathbf{X})\right),\text{ where}
\\
\scalebox{0.93}{
$\underset {\begin{subarray}{l} i=0,\ldots,\frac{d}{2}-1;\\ n\in[ 0,|\mathbf{x}|);\end{subarray}}{\text{RoPE}_(n)}= {\left[\cdots,\cos\left(\mathbb{I}(\hat{\lambda}_{i}, \hat{n})\frac{n}{\beta^i}\right),\sin\left(\mathbb{I}(\hat{\lambda}_{i}, \hat{n})\frac{n}{\beta^i}\right),\cdots\right] }$}\\
	\text{where} \quad \mathbb{I}(\hat{\lambda}_{i},\hat{n})=\begin{cases}
	1&\mbox{\quad  $n<\hat{n}$}\\
{\frac{1}{\lambda}_{i}}&\mbox{\quad $n\geq\hat{n}$}
\end{cases}
	\end{gathered}
\end{equation}

In this framework, we define a collection of rescale factors, $\mathbb{I}(\hat{\lambda}_{i},\hat{n})$, that account for two types of non-uniformity. Here, $\hat{\lambda}_i$ and $\hat{n}$ represent irregularities related to the dimensions of RoPE and the locations of tokens, respectively. More precisely, $\mathbb{I}(\hat{\lambda}_{i},\hat{n})$ is employed to adjust the rotation angle associated with the $i^{th}$ dimension of RoPE; $\hat{\lambda}_i$ functions as the specific rescaling parameter, and $\hat{n}$ marks a positional threshold for tokens. For tokens in the first $\hat{n}$-1 positions, the rescaling via $\hat{\lambda}_i$ does not apply; instead, the default RoPE rotational angle $\frac{n}{\beta^i}$ is utilized. Once tokens reach positions $n\ge\hat{n}$, the rescale factor comes into effect.

Suppose we set the desired context window to $L'$. In this scenario, our aim is to determine the optimal rescaling factors, denoted as $\mathbb{I}(\hat{\lambda}_{0},\hat{n}), \mathbb{I}(\hat{\lambda}_{1},\hat{n}), \ldots, \mathbb{I}(\hat{\lambda}_{i},\hat{n})$, for each RoPE dimension from the $1^{st}$ to the $d^{th}$. With these adjusted rescaling factors applied to the $\text{RoPE}$, the given $\text{LLM}$ should minimize the next-token prediction loss $\mathcal{L}$ (representing perplexity) on datasets $\mathbf{X}$ whose token lengths are $L'$.

\subsection{Searching the Non-uniform Position Interpolation}

To address the challenge described in Eq.~\ref{eq:problem}, we present a straightforward but powerful approach that identifies optimal RoPE rescaling coefficients, enabling comprehensive utilization of the multidimensional irregularities present in position embedding.

\textbf{Search space}. A comprehensive search space is constructed to encompass both non-uniformities. Table~\ref{tbl:searchspace} provides a detailed overview of the search space. In particular, the search process incorporates the optimization of distinct rescale factors for each dimension within the RoPE embedding. For simplification, instead of directly optimizing $\mathbb{I}(\hat{\lambda}_{i},\hat{n})$, we focus on searching for the parameters $\lambda_i$ and $\hat{n}$, where $\hat{\lambda}_{i}=1/\lambda_i$. According to Table~\ref{tbl:searchspace}, the parameter $\lambda_i$ is explored within the range from 1.0 (representing pure extrapolation) up to $s\times1.25$ (which represents interpolation with greater magnitude than PI), using increments of 0.01. Here, $s$ denotes the ratio corresponding to the desired extension of the context window.

The parameter $\hat{n}$ specifies how many leading token positions maintain their original RoPE embeddings, foregoing position interpolation. In practice, $\hat{n}$ is selected from the set \{0, 1, 2, 4, 8, 12, 16, 20, 24, 28, 32, 64, 128, 256\} through empirical evaluation. Setting $\hat{n}=0$ implies that every token position employs the optimized rescale factors.

\textbf{Evolution-based search}. The search space outlined in Table~\ref{tbl:searchspace} encompasses a wide array of positional interpolation configurations, making comprehensive exploration computationally prohibitive. For instance, an extension with $s=4\times$ results in $400^{128/2}\times14=4\times10^{167}$ possible selections. As the extension ratio increases, this search space grows at an exponential rate. To overcome these limitations, we employ an evolution search strategy~\cite{guo2020single} and incorporate two optimization methods designed to enhance the efficiency of the search process. The overall procedure is presented in Algorithm~\ref{alg:search}.

\textit{Optimized initial population generation}. Rather than relying solely on random initialization for a population of $P$ rescale factors, we explicitly include the RoPE rescale factors associated with PI, NTK, and YaRN as separate individuals at the outset. For the other $P-3$ individuals in the initial population, we generate them by introducing random mutations to these three rescale factors with a mutation probability of $p$.

\textit{Monotonically non-decreasing constraint}. Once the initial population is created, we evaluate the LLM perplexity for every candidate. To do this, the specified RoPE rescale factors are incorporated into the target LLM, after which we determine the perplexity score for input $\mathbf{X}$. The individuals with the top $k$ perplexity results are then selected as parents for the next evolutionary step. Despite this process, the enormity of the search space means that uninformed mutation and crossover operations can frequently yield suboptimal candidates, resulting in excessive, redundant perplexity evaluations. This inefficiency is especially pronounced when $L'$ is considerable, since each perplexity measurement requires extensive inference time.

To mitigate this issue, we enforce a non-decreasing monotonicity condition on the selected RoPE rescaling coefficients, such that $\lambda_i \leq \lambda_{i+1}$. Only those RoPE factors satisfying this criterion are used within the LLM during perplexity assessments, thereby considerably curtailing the search space. More precisely, our approach stipulates that $\lambda_i$ must increase consistently with respect to the RoPE dimension (namely, $i=0,\ldots,63$). The rationale for this dimensional monotonicity draws from NTK theory~\cite{jacot2018neural,tancik2020fourier,ntk}, which indicates that the dimensions associated with higher frequency (i.e., lower indices) demand less interpolation (corresponding to smaller $\lambda_i$ values), whereas those dimensions aligned with lower frequency (higher indices) can support greater interpolation, thus necessitating larger $\lambda_i$.

\begin{algorithm}[t]
	\small
	\caption{The search algorithm}
	\textbf{Input:} target LLM, input samples $\mathbf{X}$,   population size $P$, mutation size $N_1$, crossover size $N_2$, max iterations $\mathcal{T}$, mutate probability $p$\\

	\begin{algorithmic}[1]
		\label{alg:search}
		\STATE Top-k $\leftarrow$ $\phi$;
		\STATE $\text{P}_0$ $\leftarrow$ \textit{Initialize\_population\_with\_optimization} ($P$, $\mathbf{X}$, $p$);
		\FOR{$i$=1 \textbf{to} $\mathcal{T}$}
		\STATE \textit{Compute\_perplexity} (LLM, $\text{P}_{i-1}$, $\mathbf{X}$);
		\STATE  Top-k $\leftarrow$ \textit{Update\_Topk} (Top-k, $\text{P}_{i-1}$);
		\STATE$\text{P}_{mutation}$ $\leftarrow$ \textit{Mutation\_with\_mono\_constraint} (Top-k, $N_1$, $p$); 
		\STATE $\text{P}_{crossover}$ $\leftarrow$ \textit{Crossover\_with\_mono\_constraint} (Top-k, $N_2$);
		\STATE $\text{P}_i$ $\leftarrow$ $\text{P}_{mutation}$ $\cup$ $\text{P}_{crossover}$ $\cup$ Top-k;
		\ENDFOR
		\STATE Output the member with minimum perplexity value from Top-k;
	\end{algorithmic}
\end{algorithm}

\textbf{8$\times$ extension achieved without fine-tuning}. Through the use of an evolutionary search strategy, our approach discovers optimal non-uniform RoPE rescale coefficients that retain critical dimensional features and positional elements, thereby reducing information degradation due to interpolation. As shown in Fig.\ref{fig:non-ft}, this technique allows for expanding LLaMA2's context window from 4k to 32k with no requirement for fine-tuning. In comparison, methods like PI and non-uniform NTK and YaRN result in a sharp increase in perplexity following a 2$\times$ extension.

\subsection{Extending LLM Context Window to 2048K}
\label{sec:extendto2048k}

\textbf{Progressive extension to 2048k}. In this section, we present our approach for increasing the context window of pre-trained LLMs from the conventional 4k up to more than 2048k tokens. Our results show that non-uniform positional interpolation enables up to an 8$\times$ context length expansion without the need for fine-tuning. When even greater extensions are sought—such as a 512$\times$ enlargement—fine-tuning becomes essential. A potential strategy involves searching for appropriate RoPE rescaling factors tailored to the 2048k context length prior to fine-tuning. However, this approach is often hindered by the immense computational cost associated with training. Additionally, in our experience, achieving effective fine-tuning with such a significant extension ratio poses considerable difficulties (refer to Appendix for further discussion).

LongRoPE demonstrates efficacy across both baseline and fine-tuned, length-extended LLMs. To capitalize on this, we propose a resource-efficient, incremental strategy that enables reaching the desired 2048k context window using only 1k fine-tuning iterations, each limited to a training context of 256k tokens.

$\Diamond$ \textit{Scaling pre-trained LLMs to a 256k context via LongRoPE exploration}. Using LLaMA2 as a reference, we systematically investigate extending the context window to 128k and 256k tokens. This phase involves enlarging the context window by factors of 32$\times$ and 64$\times$, respectively.

$\Diamond$ \textit{Fine-tuning to 256k}. Subsequently, we adapt the pre-trained LLM to support a 256k context window through additional fine-tuning. Concretely, we begin by fine-tuning LLaMA2 for 400 steps with RoPE rescaling set for 128k. After this phase, we substitute the RoPE rescaled factors to correspond to 256k on the resulting checkpoint, followed by another 600 steps of fine-tuning. Compared to initiating fine-tuning at 256k directly, this staged process is found to be more effective.

$\Diamond$ \textit{Applying LongRoPE search to increase the context window of a fine-tuned extended LLM to 2048k}. In the final step, we carry out a second search process on the LLM that has previously been fine-tuned for a context length of 256k. Through this method, we achieve a remarkable expansion of the context window, reaching 2048k, and notably, this is accomplished without additional rounds of fine-tuning. The resulting extension factor is $512\times$.

\textbf{Recovery for shorter context windows}.  When expanding the context window to an unprecedented length of 2048k, a degradation in model performance is observed within the initial context window span. This phenomenon is documented in the literature on positional interpolation~\cite{pi}, where the embedding for positions in the original context is compressed into a restricted segment of the higher-dimensional space, thereby impairing the effectiveness of the language model. In scenarios involving a 512$\times$ extension ratio, the position representations within the initial 4k context window are densely concentrated.

To address this issue, we conduct an additional evolution search on the extended LLM, specifically tuning the RoPE rescale parameters for relatively short context lengths such as 4k and 8k. For these shorter spans, we lower the upper bound on the $\lambda$ values explored, as they require less positional interpolation. At inference time, the LLM adapts the relevant RoPE rescale factors on the fly.

 \section{Experiments}

\label{sec:exp}
\subsection{Setup}
\textbf{Evaluation Tasks and Models}. In our experiments, we integrate LongRoPE into LLaMA2-7B and Mistral-7B, assessing its efficacy across three dimensions: (1) measuring perplexity scores of LLMs with extended context when processing lengthy documents; (2) conducting a Passkey retrieval evaluation, which tests how effectively the model can locate a specific passkey embedded among large amounts of distracting information; and (3) utilizing conventional LLM benchmark tests, all constrained to a 4096-token context length.

\textbf{Fine-tuning}. For LLaMA2, we apply a linear learning rate schedule starting at 2e-5 and set the global batch size to 32. The initial stage of fine-tuning spans 400 steps using the Redpajama~\cite{redpajama} dataset, organized as segments of 128k tokens framed with BOS and EOS markers. Subsequently, utilizing the resulting checkpoint, we continue fine-tuning for an additional 600 steps to extend the context window to 256k tokens. Distributed training is performed with 8 A100 GPUs for the 128k configuration via the system described in~\cite{cube}, and for the 256k context window, training scales to 16 A100 GPUs. For Mistral, we employ a fixed learning rate of 1e-6, opting for a global batch size of 64. The protocol for both 128k and 256k models is aligned with YaRN~\cite{yarn}, executing 400 steps on Together Computer's Long-Data Collections~\cite{mistraldata}, with each sequence consisting of 16k tokens. Model training for Mistral leverages 4 A100 GPUs.

\textbf{Search}. When the target window size does not exceed 256k, the following parameters are set: $P$=64, $N_1$=$N_2$=16, $p$=0.3, $\mathcal{T}$=40, with the top 32 candidates chosen for mutation and crossover in every iteration. Perplexity evaluation is performed on 5 randomly selected samples from the PG19 validation set, each with at least the intended context length. For window sizes larger than 512k, the values for population, mutation, and crossover are all reduced by half. Perplexity assessments in this regime use 3 randomly selected validation samples from the Pile-Books3~\cite{pile} dataset.

\textbf{Baselines}. To achieve a 2048k context window, we performed fine-tuning on models initially trained with 128k and 256k context lengths. This results in LongRoPE-2048k (ft=128k) for LLaMA2 and LongRoPE-2048k (ft=256k) for Mistral. We evaluate these four models alongside the leading baselines for context window extension, which include publicly available LLMs refined via positional interpolation using PI, NTK, and YaRN approaches. Among these baselines are Together-32k~\cite{together}, Code LLaMA~\cite{codellama}, LongLoRA-full-FT-100k~\cite{longlora}, YaRN-LLaMA, and YaRN-Mistral~\cite{yarn}.

\subsection{Main Results}

\label{sec:mainresult}
\textbf{Language modeling for long sequences up to 256k tokens}. Our analysis begins with a comparison to current leading extended LLMs evaluated over sequences of up to 256k tokens. To highlight the versatility of our approach, we utilize two datasets: Proof-pile~\cite{pg19} and the PG19~\cite{pile} test sets. Perplexity measurements are obtained for multiple context sizes employing a sliding window of 256 tokens. Specifically, for PG19, evaluation covers the complete test set, which includes 100 documents. For Proof-pile, adhering to YaRN~\cite{yarn}, we randomly select 10 samples, ensuring each is at least 128k tokens in length.

Table~\ref{tbl:proof-pile} and Table~\ref{tbl:pg19} present perplexity results for LLaMA2 and Mistral models extended through various interpolation techniques on Proof-pile and PG19, respectively. Two main findings stand out: \textbf{(1)} our extended models consistently achieve lower perplexity scores as evaluation length increases from 4k up to 256k, indicating improved utilization of extended contexts. \textbf{(2)} Even when the context window is enlarged to 16$\times$ its original size—a scenario in which models typically struggle to retain performance on shorter contexts—our LongRoPE-2048k models exceed the performance of leading baselines for context lengths up to 256k.

\textbf{Long sequence language modeling beyond 2000k}. In order to assess performance on exceptionally lengthy texts, the Books3~\cite{pile} dataset is employed. For efficient evaluation, a random sample of 20 books, each with a length greater than 2048k, is selected, and a sliding window of 256k tokens is applied.

Table~\ref{tbl:books3} demonstrates that LongRoPE effectively expands the context window of LLaMA2-7B and Mistral-7B models to 2048k tokens, all while maintaining perplexity levels that are on par with or even outperform the baselines for shorter contexts ranging from 8k to 128k. Significant differences in performance are noted between the LLaMA2 and Mistral models when extended to 2048k. Specifically, Mistral shows superior results at shorter context lengths but sees perplexity rise above 7 when the context length exceeds 256k. The LLaMA2 model exhibits expected trends, with perplexity generally decreasing as the context length grows, except for minor increases at 1024k and 2048k. Additionally, LongRoPE-2048k on LLaMA2 achieves better performance when fine-tuned at 256k as opposed to 128k, which can be attributed to a lower secondary extension ratio (8$\times$ vs. 16$\times$). Conversely, Mistral achieves optimal results with a fine-tuning window of 128k tokens. This outcome is primarily because Mistral’s fine-tuning at both 128k and 256k follows YaRN's protocol using a training length of 16k, limiting the model’s ability to further scale the context window post fine-tuning.

\textbf{Passkey retrieval}. Next, we examine how context window size influences performance in generative tasks by utilizing the synthetic passkey retrieval benchmark introduced by ~\cite{passkey}. In this evaluation, the model must locate a randomly selected five-digit number, termed the passkey, embedded somewhere within an extended text. Details of the prompt construction are provided in the appendix. For assessment, we conduct 10 repetitions of the retrieval task, each time positioning the passkey at a different, randomly chosen spot uniformly sampled across the evaluation context window.

Figure~\ref{fig:passkey} presents the retrieval accuracy results alongside several baselines. The performance of existing models quickly declines to zero when sequence lengths surpass 128k. Conversely, LongRoPE-LLaMA2-2048k (ft=256k) consistently achieves high retrieval accuracy ($\ge$90\%) across the full range from 4k to 2048k, demonstrating robustness in retrieving passkeys among millions of tokens. LongRoPE-Mistral-2048k (ft=128k) maintains perfect accuracy up to 1800k, but drops to 60\% at 2048k—a trend consistent with Table~\ref{tbl:books3}, which indicates a moderate increase in perplexity at this length.

\textbf{Evaluation on standard benchmarks under the original context window.} We assess the performance of LongRoPE-2048k models within their initial context length, utilizing the Hugging Face Open LLM Leaderboard~\cite{huggingfaceboard} for both zero-shot and few-shot evaluations. Specifically, we conduct 25-shot testing on ARC-Challenge~\cite{allenai:arc}, 10-shot on HellaSwag~\cite{zellers2019hellaswag}, 5-shot on MMLU~\cite{mmlu}, and a zero-shot evaluation on TruthfulQA~\cite{lin2021truthfulqa}.

Table~\ref{tbl:commonsense} demonstrates that our models deliver similar outcomes on the initial benchmark, which was intended for shorter context lengths. Notably, they exceed the performance of the original Mistral on TruthfulQA by a margin of +0.5\%. Although LongRoPE-LLaMA2-2048k—fine-tuned using a 256k context window—exhibits marginally greater declines in performance, it still maintains acceptable results across the majority of tasks.

\subsection{Ablation Results}

\textbf{Effectiveness of the second positional interpolation}. Within our progressive extension approach, the proposed search algorithm is leveraged to apply a secondary, non-uniform positional interpolation to the fine-tuned extended LLMs. To assess its utility, we performed evaluations on the fine-tuned LLaMA2-256k model, expanding its context length to 512k, 1024k, and 2048k via PI and YaRN methods. Results in Table~\ref{tbl:secondsearch} indicate that our approach to non-uniform positional interpolation preserves perplexity across varying lengths. Conversely, both PI and YaRN display a rapid increase in perplexity as the extension ratio grows.

\textbf{Recovery performance at limited context lengths.} In order to counteract reductions in effectiveness at shorter context windows, we recalibrate the RoPE scaling parameters for LongRoPE-2048k using our search procedure. More precisely, we restrict the maximum scale factors explored by the search, thus diminishing the degree of interpolation when working with 4k and 8k contexts. Table~\ref{tbl:shortperformance} presents a perplexity comparison for LongRoPE-LLaMA2-2048k on Proof-pile at both 4k and 8k tokens, and reports average accuracy across various LLM benchmarks. The results indicate a marked enhancement in performance at these shorter context sizes.

\textbf{Investigation of the two sources of non-uniformities}. To discern the contribution of each non-uniformity type to model effectiveness, we conduct ablation studies. Our methodology involves two sets of experiments: (i) extending LLaMA2-7B to relatively short sequence lengths of 16k and 32k using three approaches—PI, optimizing solely for the RoPE dimension, and optimizing for both types of non-uniformities; (ii) taking our LLaMA2 model fine-tuned to 256k tokens and pushing it to 2048k with identical extension techniques. Perplexity metrics are gathered in the absence of further fine-tuning. As presented in Table~\ref{tbl:searchdimension}, adjustment for non-uniformity within the RoPE dimension yields a notably lower perplexity compared to the linear interpolation employed by PI. Addressing non-uniformity in token positions noticeably boosts model accuracy at 16k and 32k, yet its effects diminish at the extreme sequence length of 2048k—potentially a consequence of the vast context span. Retaining only the original token positions, without any form of interpolation, offers no tangible benefit under these circumstances; exploration of this aspect remains as prospective work.

\section{Related Works}

Beyond interpolation-based methods, this section reviews other techniques found in related research.

\textbf{Retrieval-based} approaches rely on an external memory component to store lengthy historical contexts and incorporate retrieval mechanisms that extract relevant information during inference~\cite{longllama,wang2023augmenting,borgeaud2022improving}. These strategies often involve substantial alterations to LLM architectures. By comparison, our approach is considerably more lightweight, only requiring minimal changes to positional embeddings. Additionally, we address a wider array of long-context tasks outside the scope of retrieval, such as comprehensive document summarization and few-shot learning scenarios.

\textbf{Attention-driven expansion of context windows}. In addition to the interpolation of positional embeddings, various studies have managed to expand the effective input context window while retaining the original length of an LLM’s context window by adapting the underlying attention mechanism~\cite{han2023lminfinite, streamingllm,parallelwindow}. The principal strategy involves curbing the excessive growth in attention computation resulting from new positional tokens by leveraging innovative attention mask designs. The techniques in this category can be combined with positional interpolation approaches to offer complementary benefits.

\textbf{Fine-tuning based approaches} investigate strategies to optimize the adaptation of pre-trained LLMs with adjusted position embeddings to support longer context windows. Several studies, such as Code LLaMA~\cite{codellama}, LLaMA2 Long~\cite{xiong2023effective}, and ScaledRoPE~\cite{liu2023scaling}, employ substantially increased base values for RoPE and conduct fine-tuning at the desired context lengths. Our approach provides adaptability to multiple target lengths and demonstrates capability with contexts surpassing 2M tokens. In response to the significant GPU requirements imposed by fine-tuning for extended context lengths (e.g., exceeding 128k), recent works including LongLoRA~\cite{longlora} and PoSE~\cite{zhu2023pose} have been introduced to alleviate computational overhead. Notably, our method is complementary to these resource-efficient fine-tuning techniques.

\section{Conclusion}

This paper introduces LongRoPE, a novel technique that dramatically pushes the context window size of LLMs to 2048k tokens, all while preserving performance across shorter, conventional windows. Our approach leverages two distinct non-uniform characteristics in RoPE positional embedding, harnessed through an efficient evolutionary search algorithm. This confers dual advantages: first, it delivers robust initialization for subsequent fine-tuning, and second, it allows an 8$\times$ increase in context window without the need for any fine-tuning. Building on this foundation, we put forth a gradual extension framework, wherein 256k-token fine-tuned LLMs are incrementally expanded to support a 2048k-token window—achieved without supplementary fine-tuning. Through comprehensive experimentation, we demonstrate LongRoPE’s strong efficacy. We anticipate that the LongRoPE-2048k model family will facilitate novel applications requiring extended contexts and act as a catalyst for continued advancements in this area.

\section*{Broader Impacts} The aim of this paper is to contribute to the progression of the Machine Learning discipline. While the research could potentially have a range of societal effects, we have not identified any particular impacts that require explicit discussion at this time.

	\balance

\end{document}